\section{Acid, base và pH}

\subsection{Acid và base}

Trong nước, acid tạo ra ion \ce{H+}; base tạo ra ion \ce{OH-}. Những hydroxide tan tốt trong nước được gọi là \textbf{kiềm}.

\begin{ghinho}{Nhận biết nhanh}
\begin{itemize}
\item Dung dịch acid: làm quỳ tím hóa đỏ.
\item Dung dịch kiềm: làm quỳ tím hóa xanh; phenolphthalein chuyển màu hồng.
\item Base không tan thường được nhận biết qua chất rắn kết tủa, chẳng hạn \ce{Cu(OH)2}, \ce{Fe(OH)3}.
\end{itemize}
\end{ghinho}

Một số acid thường gặp: \ce{HCl}, \ce{H2SO4}, \ce{CH3COOH}. Một số base thường gặp: \ce{NaOH}, \ce{KOH}, \ce{Ca(OH)2}, \ce{Cu(OH)2}.

\subsection{Tính chất cơ bản của acid}

Các phản ứng thường gặp ở mức THCS:

\begin{enumerate}
\item Acid làm đổi màu chất chỉ thị.
\item Acid phản ứng với một số kim loại tạo muối và khí hydrogen.
\item Acid phản ứng với base tạo muối và nước.
\item Acid phản ứng với oxide base tạo muối và nước.
\item Acid phản ứng với một số muối; phản ứng thường xảy ra khi tạo kết tủa, khí hoặc chất điện li yếu như nước.
\end{enumerate}

Ví dụ:
\[
\ce{Zn + 2HCl -> ZnCl2 + H2 ^},
\]
\[
\ce{HCl + NaOH -> NaCl + H2O},
\]
\[
\ce{CuO + 2HCl -> CuCl2 + H2O},
\]
\[
\ce{Na2CO3 + 2HCl -> 2NaCl + CO2 ^ + H2O}.
\]

\subsection{Tính chất cơ bản của base}

\begin{enumerate}
\item Dung dịch kiềm làm đổi màu chất chỉ thị.
\item Base phản ứng với acid tạo muối và nước.
\item Dung dịch kiềm có thể phản ứng với oxide acid.
\item Dung dịch kiềm có thể phản ứng với một số dung dịch muối tạo hydroxide không tan.
\end{enumerate}

Ví dụ:
\[
\ce{2NaOH + CO2 -> Na2CO3 + H2O},
\]
\[
\ce{CuSO4 + 2NaOH -> Cu(OH)2 v + Na2SO4}.
\]

\begin{notechem}
Khi viết phản ứng, không nên học thuộc từng phương trình riêng lẻ. Hãy xác định loại chất ở hai vế, dự đoán sản phẩm rồi kiểm tra điều kiện để phản ứng xảy ra.
\end{notechem}

\subsection{Thang pH}

Thang pH được dùng để đánh giá môi trường của dung dịch:
\[
pH<7:\ \text{môi trường acid};\qquad
pH=7:\ \text{trung tính};\qquad
pH>7:\ \text{môi trường base}.
\]

Ở chuyên đề này, pH chủ yếu được dùng để \textbf{nhận xét môi trường và giải thích tình huống thực tế}. Không cần dùng công thức logarithm để tính pH.

\begin{vidu}
Một mẫu đất có pH giảm sau thời gian dài sử dụng phân đạm. Môi trường đất thay đổi như thế nào?
\end{vidu}
\begin{loigiai}
Khi pH giảm, môi trường đất trở nên acid hơn.

\textbf{Vậy,} đất có xu hướng chua hơn.
\end{loigiai}

\begin{tuluyen}
Cho các dung dịch có pH lần lượt bằng 3, 7 và 10. Xác định môi trường của mỗi dung dịch.
\dapso{pH 3: acid; pH 7: trung tính; pH 10: base.}
\begin{loigiaiBT}
pH nhỏ hơn 7 ứng với môi trường acid; pH bằng 7 ứng với môi trường trung tính; pH lớn hơn 7 ứng với môi trường base.

\textbf{Vậy,} ba dung dịch lần lượt có môi trường acid, trung tính và base.
\end{loigiaiBT}
\end{tuluyen}

\begin{tuluyen}
Nêu hiện tượng và viết phương trình hóa học khi cho một mẩu zinc vào dung dịch hydrochloric acid loãng.
\dapso{Zinc tan dần, có khí không màu thoát ra; \ce{Zn + 2HCl -> ZnCl2 + H2 ^}.}
\begin{loigiaiBT}
Zinc phản ứng với hydrochloric acid tạo zinc chloride và hydrogen:
\[
\ce{Zn + 2HCl -> ZnCl2 + H2 ^}.
\]
Zinc tan dần và xuất hiện bọt khí không màu.

\textbf{Vậy,} hiện tượng là zinc tan dần và có khí hydrogen thoát ra.
\end{loigiaiBT}
\end{tuluyen}
